\section{Original Madgwick recurrence as an exact operator graph}
\label{sec:exact-madgwick}

\begin{panel}
\textbf{Status.} This chapter defines a mathematical profile and gives conditional
proofs. It does not supply a new runtime, machine-checked proof, simulation,
benchmark or physical accuracy result. Its purpose is to retain the chosen
operators and their dependencies before a rounded numerical state is materialized.
The verified R10 runtime remains a separate preserved implementation.
\end{panel}

\subsection{Source identity and discrete profile}
The source is Sebastian O. H. Madgwick's 2010 report,
\href{https://x-io.co.uk/downloads/madgwick_internal_report.pdf}{\emph{An efficient
orientation filter for inertial and inertial/magnetic sensor arrays}}.
We select its simplified recurrence, equations (25)--(34), (42)--(51), under the
name \texttt{MDG2010-EQ-SIMPLIFIED-R1}. Normalization, sample ordering and
otherwise undefined branches are specified below. The finite-mixture family
from equations (35)--(38) is retained as a distinct alternative, not an identity
used to rewrite the simplified recurrence.

The current \href{https://github.com/xioTechnologies/Fusion}{x-io Fusion library}
describes a revised thesis algorithm. It is not silently substituted for the
selected 2010 profile. Here ``fusion'' in the finite-mixture equations names the
source's combination of estimates, not that later software package.

Equation (45) forms the magnetic reference from the previous attitude and the
current normalized magnetic sample. The report's appendix C code instead stores
a magnetic reference after its new attitude update for use at the next sample.
It also prepares \texttt{twom\_x/y/z} before normalizing the magnetic input and
later uses those aliases for the flux calculation. These orders have distinct
literal meanings. This chapter pins the equation profile; it asserts neither
bit identity with that C code nor equivalence to its unnormalized aliases.

\subsection{Typed state, units and frame convention}
Let $B$ be the calibrated right-handed sensor/body frame and $E$ the right-handed
filter reference frame. The positive $E$-$z$ axis is the declared accelerometer
reference direction after sensor-sign calibration; its horizontal $x$ axis is
the chosen magnetic or heading reference. The alignment of $E$ with station ENU
is an explicit transform, not an assumed identity.

The scalar-first Hamilton quaternion
\[
q_{\rm att}=q=(w,x,y,z)=(w,\boldsymbol v),\qquad q^Tq=1
\]
maps body vectors to $E$ vectors. The separate ASA/JK word is denoted
$q_{\rm word}\in\{0,\ldots,2^{32}-1\}$. The quaternion and Boolean word have
different types and no shared state fields. For quaternions $a=(a_0,\boldsymbol a)$
and $b=(b_0,\boldsymbol b)$, define
\[
a\otimes b=(a_0b_0-\boldsymbol a\cdot\boldsymbol b,
 a_0\boldsymbol b+b_0\boldsymbol a+\boldsymbol a\times\boldsymbol b),
\]
\[
q^*=(w,-\boldsymbol v),\qquad
[0,\boldsymbol p_E]=q\otimes[0,\boldsymbol p_B]\otimes q^*.
\]
All subsequent frame equations follow this convention directly.

The ordered event is
\[
I_n=(\mathrm{id}_n,t_n,\boldsymbol\omega_n,\boldsymbol a_n,
 \boldsymbol m_n,\mathrm{validity}_n,\mathrm{calibrationId}),
\qquad h_n=t_n-t_{n-1}>0.
\]
Time is in seconds on a declared monotonic scale. Gyro samples are calibrated
rad/s; accelerometer values are m/s$^2$ before direction normalization; magnetic
components have one declared consistent unit. Accelerometer specific force must
receive the declared sign convention before it is interpreted as the reference
direction. A clock mapping to GPST is required for station/orbit composition.

Persistent data comprise the quaternion, previous event time/index, and optional
gyro bias $\boldsymbol b_\omega$ in rad/s. Parameters include $\beta\ge0$ in
s$^{-1}$, $\zeta\ge0$ in s$^{-2}$, fixed calibration, initialization, reference
policy, clock/frame/mounting transforms and domain branches. A lagged magnetic
variant additionally stores $(b_x,b_z)$ and its initialization; the selected
current-reference variant recomputes it and needs no independent persistent
magnetic state. Fixed calibration bias and estimated drift are distinct terms.

Raw integer ADC values with exact scale ratios are rational literals. A finite
IEEE value has an exact dyadic-rational meaning; an intended decimal ratio has
a different literal type. Raw bit provenance can remain distinct from either
mathematical value. Exact degree-to-radian conversion introduces symbolic $\pi$.
Thus the purely algebraic profile assumes rational/algebraic radian inputs;
symbolic $\pi$ conversion belongs to an explicitly broader operator domain.
An exact recorded value does not certify the corresponding physical measurement.

\subsection{Normalization, rotation and residual operators}
For a real vector $v\ne0$, define the partial operator
\[
\mathcal N(v)=\frac{v}{\sqrt{v^Tv}},
\]
where the square root is the positive root. Admitted observations give
$\widehat{\boldsymbol a}=\mathcal N(\boldsymbol a)$ and
$\widehat{\boldsymbol m}=\mathcal N(\boldsymbol m)$.
For unit $q$, the body-to-reference rotation is
\[
R(q)=\begin{bmatrix}
1-2(y^2+z^2)&2(xy-wz)&2(xz+wy)\\
2(xy+wz)&1-2(x^2+z^2)&2(yz-wx)\\
2(xz-wy)&2(yz+wx)&1-2(x^2+y^2)
\end{bmatrix}.
\]
Substitution into $R(q)^T(0,0,1)^T-\widehat{\boldsymbol a}$ yields the
accelerometer objective and its ambient Jacobian:
\begin{align}
f_g(q)&=\begin{bmatrix}
2(xz-wy)-\widehat a_x\\2(wx+yz)-\widehat a_y\\
1-2(x^2+y^2)-\widehat a_z
\end{bmatrix},\label{eq:mad-fg}\\
J_g(q)&=\begin{bmatrix}
-2y&2z&-2w&2x\\2x&2w&2z&2y\\0&-4x&-4y&0
\end{bmatrix}.\label{eq:mad-jg}
\end{align}
For a fixed reference $\boldsymbol b_E=(b_x,0,b_z)^T$, the magnetic residual is
\begin{equation}
f_b(q)=\begin{bmatrix}
b_x(1-2(y^2+z^2))+2b_z(xz-wy)-\widehat m_x\\
2b_x(xy-wz)+2b_z(wx+yz)-\widehat m_y\\
2b_x(wy+xz)+b_z(1-2(x^2+y^2))-\widehat m_z
\end{bmatrix}.
\label{eq:mad-fb}
\end{equation}
Its Jacobian, with columns ordered $(w,x,y,z)$, is
\begin{equation}
\begin{split}
J_b(q)=\bigl[\;&c_w\quad c_x\quad c_y\quad c_z\;\bigr],\\
c_w&=\begin{bmatrix}-2b_zy\\-2b_xz+2b_zx\\2b_xy\end{bmatrix},\qquad
c_x=\begin{bmatrix}2b_zz\\2b_xy+2b_zw\\2b_xz-4b_zx\end{bmatrix},\\
c_y&=\begin{bmatrix}-4b_xy-2b_zw\\2b_xx+2b_zz\\2b_xw-4b_zy\end{bmatrix},\qquad
c_z=\begin{bmatrix}-4b_xz+2b_zx\\-2b_xw+2b_zy\\2b_xx\end{bmatrix}.
\end{split}
\label{eq:mad-jb}
\end{equation}
Products such as $b_xx$ mean $b_x\,x$, not an extra subscript. These expressions
follow by direct polynomial differentiation: for example,
$\partial f_{b,1}/\partial y=-4b_xy-2b_zw$ and
$\partial f_{g,3}/\partial w=0$.

For IMU, set $f=f_g$, $J=J_g$. For MARG, stack residuals and Jacobians vertically:
\[
f=\begin{bmatrix}f_g\\f_b\end{bmatrix},\quad
J=\begin{bmatrix}J_g\\J_b\end{bmatrix},\quad
F(q)=\tfrac12 f(q)^Tf(q),\quad
s=\nabla F=J^Tf.
\]
The chain rule gives $s=J_g^Tf_g+J_b^Tf_b$ in MARG mode. These are the source's
unweighted residuals. Changing sensor weights defines a different recurrence.
The derivatives refer to the displayed ambient polynomial extension evaluated
at a unit quaternion. Agreement of two residual formulas on the unit sphere
does not imply equality of their ambient gradients.

\subsection{Magnetic reference and optional drift state}
Before the current MARG residual is formed, compute
\begin{align}
\boldsymbol h_n&=R(q_{n-1})\widehat{\boldsymbol m}_n,\\
b_{x,n}&=\sqrt{h_{x,n}^2+h_{y,n}^2},&b_{z,n}&=h_{z,n}.
\label{eq:mad-reference}
\end{align}
Hold these values fixed while evaluating $J_b$. Differentiating through the
reference update would add chain-rule terms absent from the chosen source
gradient. A dependency in the event graph therefore need not be a differentiation
variable of every local objective.

\begin{proposition}[Reference norm]
If $q_{n-1}$ and $\widehat{\boldsymbol m}_n$ have unit norm, then
$b_{x,n}\ge0$ and $b_{x,n}^2+b_{z,n}^2=1$.
\end{proposition}
\begin{proof}
Quaternion conjugation by a unit quaternion preserves vector norm. The square
of $b_x$ is $h_x^2+h_y^2$, so adding $b_z^2$ gives $\boldsymbol h^T\boldsymbol h=1$.
This proves a representation identity, not an undisturbed magnetic measurement.
\end{proof}

Define the correction direction by an explicit zero branch:
\[
e=\begin{cases}s/\sqrt{s^Ts},&s^Ts>0,\\0,&s^Ts=0.\end{cases}
\]
Zero gradient means no gradient correction, not proof that the attitude is true.
For enabled MARG drift with both observations admitted and $b_x>0$, use
\begin{align}
\boldsymbol d_\omega&=2\operatorname{vec}(q_{n-1}^*\otimes e),\\
\boldsymbol b_{\omega,n}&=\boldsymbol b_{\omega,n-1}
 +\zeta h_n\boldsymbol d_\omega,&
\boldsymbol\omega_c&=\boldsymbol\omega_n-\boldsymbol b_{\omega,n}.
\label{eq:mad-bias}
\end{align}
The updated bias is used in this sample. Otherwise the bias is held. For
$e=(e_w,e_x,e_y,e_z)$, the drift direction is explicitly
\[
\boldsymbol d_\omega=2\begin{bmatrix}
we_x-xe_w-ye_z+ze_y\\
we_y+xe_z-ye_w-ze_x\\
we_z-xe_y+ye_x-ze_w
\end{bmatrix}.
\]
The vector projection discards the scalar quaternion component. The ambient
gradient need not be tangent to the unit sphere, so that scalar component must
not silently be interpreted as a gyro axis. The selected IMU mode does not
adapt all-axis gyro bias from an unobservable heading.

The source's gain parameterization can remain symbolic:
\[
\beta=\frac{\sqrt3}{2}\,\widetilde\omega_\beta,\qquad
\zeta=\frac{\sqrt3}{2}\,\widetilde{\dot\omega}_\zeta.
\]
The error and drift-rate quantities are declared parameters; these identities
do not establish their physical values or prove filter convergence.

\subsection{The selected update and its unit-norm invariant}
With corrected body rate $\boldsymbol\omega_c$, define
\begin{align}
\dot q_\omega&=\tfrac12q_{n-1}\otimes[0,\boldsymbol\omega_c],\\
u_n&=q_{n-1}+h_n(\dot q_\omega-\beta e),\\
q_n&=\mathcal N(u_n),\qquad u_n^Tu_n>0.
\label{eq:mad-step}
\end{align}
All candidate state fields, including bias, are committed together only after
the candidate quaternion satisfies the normalization domain.

\begin{proposition}[Unit norm after every successful update]
Starting with a unit quaternion, any finite sequence of successful updates
\eqref{eq:mad-step} produces unit quaternions.
\end{proposition}
\begin{proof}
For $S=u_n^Tu_n>0$, the positive-root semantics give
$\mathcal N(u_n)^T\mathcal N(u_n)=S/(\sqrt S)^2=1$.
Induction over the accepted events proves the statement. Existence of a valid
candidate and correct evaluation of the declared branches remain premises.
\end{proof}

To preserve the source's other derivation without conflating algorithms, let
\[
\mu=\alpha\|\dot q_\omega\|h,\quad \alpha>1,\qquad
\gamma=\frac{\beta}{\mu/h+\beta}
\]
when the denominator is positive. The finite-mixture candidate is
\begin{align*}
u_{\rm mix}&=q+(1-\gamma)h\dot q_\omega-\gamma\mu e\\
 &=q+(1-\gamma)h(\dot q_\omega-\beta e),
\end{align*}
because $\gamma\mu=(1-\gamma)\beta h$. If both denominator terms vanish, the
declared extension chooses $\gamma=0$.
At finite $\gamma$, this is generally a different effective step from
\eqref{eq:mad-step}. In particular, $\dot q_\omega=0$ and $\beta>0$ give
$\mu=0$, $\gamma=1$, and $u_{\rm mix}=q$, even if $e\ne0$. The large-$\alpha$
simplification is not an identity at every stationary input. An exact graph
preserves the chosen recurrence; it does not remove this model choice.

\subsection{Domain branches are part of the formal input}
The following are explicit definitions resolving partial equations; they are
not presented as pre-existing empirical rejection heuristics.
\begin{center}\begin{tabularx}{\linewidth}{@{}p{.34\linewidth}X@{}}\toprule
Condition & Declared action\\\midrule
Zero initial quaternion & No initialized state; require a declared nonzero initialization.\\
Untyped/nonfinite input, missing gyro or invalid clock mapping & No transition; record the input-domain failure.\\
Duplicate or decreasing event time & No transition; a separate input policy must resolve ordering or duplicates.\\
Zero or missing acceleration & Gyro-only update; no magnetic correction or drift adaptation.\\
Acceleration admitted; magnetic sample zero or missing & IMU objective; hold all-axis drift estimate.\\
Both vectors admitted and $b_x=0$ & MARG residual remains defined; mark heading degeneracy and hold all-axis drift adaptation.\\
Exactly zero gradient & Set $e=0$; no gradient-driven bias increment.\\
$\beta=0$ & No attitude-gradient correction; independently enabled $\zeta$ remains a separate operation.\\
Exactly zero candidate $u_n$ & Hold the last complete state and report undefined normalization.\\
Observation-quality flag & Use the declared input flag; do not invent an undeclared small-norm epsilon or outlier rule.\\\bottomrule
\end{tabularx}\end{center}
An algebraic domain supports exact sign/equality decisions in principle; the cost
is not bounded here. A broader transcendental domain must retain unresolved
comparisons rather than silently treating a finite-precision threshold as exact.

\subsection{Exact graph representation and conditional rewrites}
The graph consists of typed literals, addition, subtraction, multiplication,
division, positive square root, tuples/projections, dot/cross/Hamilton products
and explicit branches. Roots bind the state to its operator profile, samples,
clock and calibration. Division nodes carry nonzero-denominator obligations;
square-root nodes carry nonnegative-operand obligations. With algebraic inputs,
closure under these operators gives algebraic outputs for every defined finite
event prefix. It does not give a finite machine-word bound on their representation.

\begin{proposition}[Positive-scale cancellation]
For $c>0$ and $v\ne0$, $\mathcal N(cv)=\mathcal N(v)$.
\end{proposition}
\begin{proof}
$\sqrt{c^2v^Tv}=c\sqrt{v^Tv}$ by the positive-root convention; the factors cancel.
For $c<0$ the normalized representative changes sign, and for $c=0$ it is undefined.
\end{proof}

\begin{proposition}[Rotation consumer without a normalization square root]
For quaternion $v\ne0$, $S=v^Tv$ and pure vector quaternion $[0,p]$,
\[
R(\mathcal N(v))p=\frac{\operatorname{vec}(v\otimes[0,p]\otimes v^*)}{S}.
\]
\end{proposition}
\begin{proof}
Each quaternion factor contributes $1/\sqrt S$. Their product contributes
$1/S$. This removes the square root for this rotation consumer only.
\end{proof}

\begin{proposition}[Pure-gyro projective step]
If the corrected rate is fixed for a step and there is no attitude-gradient
correction, then, wherever defined,
\[
\mathcal N\!\left(\mathcal N(v)+\tfrac h2\mathcal N(v)
 \otimes[0,\boldsymbol\omega_c]\right)
=\mathcal N\!\left(v+\tfrac h2v\otimes[0,\boldsymbol\omega_c]\right).
\]
\end{proposition}
\begin{proof}
Factor the positive common factor $1/\sqrt{v^Tv}$ from the candidate and apply
positive-scale cancellation. The rate must still be obtained under its declared
attitude/bias convention, not recomputed using an unnormalized substitution.
For real $h$ and corrected rate, the candidate is automatically nonzero:
Hamilton norm multiplicativity gives
\[
 \left\|v+\tfrac h2v\otimes[0,\boldsymbol\omega_c]\right\|^2
 =(v^Tv)\left(1+\tfrac{h^2}{4}\|\boldsymbol\omega_c\|^2\right)>0.
\]
This establishes the normalization domain for this subcase only.
\end{proof}

The corrected homogeneous numerator is instead
\[
v+\tfrac h2v\otimes[0,\boldsymbol\omega_c]
 -h\beta\sqrt{v^Tv}\,e(\mathcal N(v)).
\]
Deleting its square root or substituting $v$ into the unit-quaternion source
polynomials changes the recurrence in general. Moreover, differentiation of the
composed normalized objective gives
\[
\nabla_v F(\mathcal N(v))
=\frac{(I-qq^T)s}{\|v\|},\qquad q=\mathcal N(v).
\]
This is a projected gradient, not the source's ambient $s$. Therefore a
projective optimization may apply the proved local identities above, but may
not silently replace the descent direction or its effective magnitude.

\begin{proposition}[Quaternion sign cover]
The selected recurrence is sign-equivariant wherever defined: replacing its
initial quaternion by $-q$ changes every quaternion representative's sign while
leaving its represented rotation and gyro-bias increments unchanged.
\end{proposition}
\begin{proof}
$R(-q)=R(q)$ because rotation is quadratic. The residuals are even in $q$,
their displayed Jacobians are odd, and thus $e(-q)=-e(q)$, including the zero
branch. The magnetic reference is unchanged. In
$q^*\otimes e$ the two signs cancel, so bias increments agree. The gyro derivative
and candidate both change sign; normalization preserves that sign relation.
Induction gives the result.
\end{proof}
This is equivalence of rotations under $q\sim-q$, not equality of stored tuples,
raw IEEE bit patterns or hashes. No automatic sign canonicalization is assumed.
Euler-angle display and shortest-arc comparison need their own explicit branch
conventions and potentially transcendental operators.

\subsection{Causal sample processing and graph cardinality}
An event reads only the preceding state, its current sample and fixed parameters.
Induction therefore proves that the state after event $n$ depends on the seed
and events $1,\ldots,n$, never on an unobserved future sample. Checkpoint restart
requires the complete selected state: attitude, optional bias, event time/index,
and any lagged reference, calibration or scheduler state that the profile retains.
Four or nine state roots do not imply four or nine fixed-size machine words.

For an input alphabet of $K$ possible records there are $K^n$ streams of length
$n$. An injective $B$-bit representation of every such stream is impossible when
$K^n>2^B$. This concerns complete stream reconstruction; it does not claim that
every distinct stream yields a distinct final quaternion. Graph sharing reduces
duplicate expression trees but cannot repeal the counting bound or guarantee
bounded algebraic degree, exact-predicate cost or evaluation memory.

A rounded checkpoint replaces exact lineage by a different numeric state. An
exact checkpoint must retain a sufficient exact representation and every
dependency not removed by a proved identity. Numerical evaluation may be deferred
to an output request; a rounded display is not fed back as the next exact state.

JK/ASA may schedule observations or exports of an attitude already updated from
all accepted sensor samples. Skipping samples, changing $h$, changing gains or
reordering measurements changes the recurrence and must be a named profile or
recorded control event. Two noncommuting angular increments already show why
two ordered quaternion steps cannot generally be replaced by one summed-rate
step. Observation cadence and sensor-integration cadence are separate contracts.

\subsection{Station composition, observability and error meaning}
Let $L$ be station ENU, $A$ the antenna/camera frame, $C_{L\leftarrow E}$ the
declared filter-to-ENU alignment, and $C_{A\leftarrow B}$ the mounting rotation.
For the simultaneous station-relative line of sight $\ell_L(t)$,
\[
\ell_A(t)=C_{A\leftarrow B}\,R(q(t))^T\,
 C_{L\leftarrow E}^T\ell_L(t).
\]
The orbital kernel supplies satellite/station geometry; attitude changes its
coordinates, not the satellite ephemeris. A moving station additionally requires
its position and velocity at the same time. An attitude quaternion supplies
neither. Between-sample interpolation, sensor latency, magnetic declination,
reference-frame alignment and mounting updates are additional declared operators.

One measured reference direction leaves rotation about that direction
unobservable. Two consistent nonparallel directions determine a rotation,
with its two quaternion representatives; $b_x=0$ makes the reference directions
collinear. Linear acceleration can violate the accelerometer reference model,
and disturbed magnetism does not become true heading through exact arithmetic.
A GNSS/INS position solution requires additional measurements and an estimator
with position/velocity/bias state. This attitude formalization establishes no
millimetre positioning or arcsecond pointing result.

Exact arithmetic of a discrete rule remains distinct from time discretization.
For a sufficiently smooth ideal derivative $F$ on a fixed branch,
\[
q(t+h)=q(t)+hF(q(t),t)+O(h^2).
\]
Evaluating the first two terms exactly does not remove the remaining terms.
Normalization is smooth away from zero, but near zero candidates, zero gradients
or mode switches, this local smoothness argument needs additional premises.
The same separation applies to R10's RK4 rule, approximated forcing/frame
coefficients, initial-state knowledge and physical force models. Sensor
calibration, quantization and noise are further independent limitations.

\begin{panel}
\textbf{Logical strategy.} Retain exact typed literals, the chosen recurrence,
sample provenance and state roots. Simplify only under proved domain and sign
conditions. Defer finite-precision materialization to a declared output boundary.
This supplies a formal route to removing intermediate arithmetic roundoff; it
does not by itself bound expression growth, discretization error or physical error.
\end{panel}
